\documentclass[12pt]{jlreq}
\usepackage{amsmath, amssymb}

\newcommand{\C}{\mathbb{C}}
\newcommand{\R}{\mathbb{R}}
\newcommand{\Z}{\mathbb{Z}}
\newcommand{\N}{\mathbb{N}}
\newcommand{\F}{\mathbb{F}}
\newcommand{\Prob}{\mathbb{P}}
\newcommand{\Exp}{\mathbb{E}}
\newcommand{\dd}{\,d}
\newcommand{\Ker}{\operatorname{Ker}}
\newcommand{\Impart}{\operatorname{Im}}
\newcommand{\Hom}{\operatorname{Hom}}

\begin{document}

\(S = k[t]\) を体 \(k\) 上の \(1\) 変数多項式環，\(K\) を \(S\) の商体とする．\(S\) の部分環 \(R\) を次のように定める：
\[
  R = k[t^4, t^{10}, t^{13}].
\]

(1) \(I = \{ x \in K \mid xS \subset R \}\) とおいたとき，\(I\) は \(R\) と \(S\)，両方のイデアルであることを示せ．

(2) \(I\) の \(R\) のイデアルとしての生成系のうち，生成元の個数が最小のものを \(1\) つ求めよ．

\end{document}